\scanpage{22}
\mspage{19}

\deleted{
Exemple : Considérons un nb réel irrationnel $r$, et considérons l'ens
\[
L=\underbrace{\mathbb Q\cap]-\infty,r]}_{L_1}
\amalg
\underbrace{(\mathbb Q\cap[r,+\infty])\times\varepsilon}_{L_2}
\]
où $\varepsilon$ est un ens \deleted{à deux} quelconque de card $\geq2$. On \deleted{peut} a donc
\diagram{\[\begin{array}{ccc@{\quad}l}
L&\overset{p}{\unsure{-|\longrightarrow}}&\mathbb Q&\text{(surjectif)}\\
L&\xrightarrow{q}&\varepsilon\amalg\{0\}&q(x)=0\text{ si }x\in L
\end{array}\]}
et si $x,y\in L$, on pose
\[
x<y\quad\Longleftrightarrow\quad px<py\quad\uline{\text{et}}
\]
}

\uline{Ex} Soit $L_0$ un ordonné \deleted{\unsure{str.}} \inserted{non vide} \deleted{filtrant croissant et filtrant} \deleted{str. filtrant} loc. str. filtrant croiss et décroissant, et telle que
\[
\exists x\in L_0,\ \exists x'<x<x'',
\]
et soit $L_i$ $(i\in I)$ une famille d'ens ordonnés satisfaisant les $=$s conditions. Soit
\[
L=L_0\amalg\coprod_iL_i.
\]
On met la relation d'ordre
\[
x<y\quad\unsure{\text{ssi}|\text{ssc}}\quad
\left|
\begin{array}{l}
\text{ou bien }x,y\text{ dans }\deleted{L_0}\inserted{\text{un }=L_0\text{ ou}}\\
\quad\text{un }L_i,\text{ et }x<y\text{ dans }\inserted{\text{cet}}L_i\\
\text{ou bien }x\in L_0,\ y\in L_i\quad(i\in I)
\end{array}
\right.
\]
Elle satisfait les $=$s conditions que $L_0,L_i$ mais n'est pas filtrante croissante !

\mspage{21}

\rotatebox{90}{\unsure{EXTRA STRONG}}
\rotatebox{90}{\unsure{EXTRA STRONG}}
